% =============================================================================
% ABSTRACT  (roman page i)
% =============================================================================
\chapter*{ABSTRACT}
\addcontentsline{toc}{chapter}{Abstract}
{\centering\fontsize{14}{18}\bfseries\underline{ABSTRACT}\par}
\vspace{0.8cm}

Inventory mismanagement is one of the most persistent operational challenges confronting small-scale
retail outlets in India, commonly referred to as kirana stores. These establishments, numbering
approximately twelve million across the country, continue to rely on manual record-keeping methods
such as physical ledgers or basic spreadsheet tools, which are ill-suited for real-time stock
monitoring, demand forecasting, or supplier evaluation. Enterprise-grade solutions designed for
large retail chains are financially prohibitive and operationally complex for single-owner
businesses that typically lack dedicated information technology personnel.

In response to this gap, a system named StockSense AI was designed and developed as an
intelligent, low-cost inventory management solution tailored specifically to the operational
context of small Indian retailers. The proposed system integrates automated workflow orchestration,
artificial intelligence-driven decision support, and real-time communication into a cohesive
platform that requires minimal technical expertise to operate.

The architecture of StockSense AI consists of four principal layers: a React~19 single-page
application serving as the user-facing frontend, an n8n-based workflow automation engine
comprising ten modular workflows constituting the backend intelligence pipeline, a Supabase
PostgreSQL database providing persistent storage and real-time data propagation, and an external
services layer incorporating the Groq Llama~3.1~8B large language model, Google Gemini~2.5
Flash multimodal model, and Twilio WhatsApp Business API.

The system automates stockout date prediction through a sales-velocity formula, classifies
products into risk tiers, generates natural-language reorder suggestions with supplier-level
reasoning, scores suppliers using a multi-criteria Weighted Sum Model, and delivers formatted
stock alerts directly to the store owner's WhatsApp account. An AI vision module enables
product onboarding from a photograph, while a separate invoice scanning workflow automates
stock-level updates from supplier delivery bills.

The prototype was evaluated against a dataset of ten products and four suppliers across multiple
inventory states. The complete intelligence pipeline executes in approximately thirteen seconds,
WhatsApp message delivery is achieved within five seconds of pipeline completion, and the
total infrastructure cost for production deployment remains below five hundred rupees per month.
The system thus demonstrates that advanced inventory intelligence, previously accessible only to
large retail enterprises, can be delivered to kirana store owners at a practically negligible cost.

\vspace{0.8cm}
\noindent\textbf{Keywords:} Inventory Management, Supply Chain Intelligence, Workflow Automation,
n8n, Large Language Model, Supplier Scoring, Kirana Stores, SME Technology, Stockout Prediction,
WhatsApp Alerts.

\clearpage
